% Source: https://artofproblemsolving.com/community/c3h2416978_1200th_post_dedicated_to_bloggy
% Stable identity: aops-topic-2416978-q18
From $B_3=0$ we need $3+p=0$, so $p=-3$. Then $B_0=\dfrac{-3}{q}=-1$ gives $q=3$. So the only candidate is
\[B_n=\dfrac{n-3}{n+3},\]
and it fits all seven terms:
\[\dfrac{-3}{3}=-1,\quad\dfrac{-2}{4}=-\dfrac12,\quad\dfrac{-1}{5},\quad\dfrac06=0,\quad\dfrac17,\quad\dfrac28=\dfrac14,\quad\dfrac39=\dfrac13.\]
(The pattern is hidden by the reduction of $-\dfrac24$, $\dfrac28$ and $\dfrac39$: the numerators go up by $1$ from $-3$ and the denominators by $1$ from $3$.) Therefore
\[B_7=\dfrac{4}{10}=\dfrac25,\qquad B_8=\dfrac{5}{11},\qquad B_9=\dfrac{6}{12}=\dfrac12,\qquad B_{10}=\dfrac{7}{13},\]
that is, $\boxed{B_7=\dfrac25,\ B_8=\dfrac5{11},\ B_9=\dfrac12,\ B_{10}=\dfrac7{13}}$.
